% !TEX root = ../diplomarbeit.tex

\section{\emph{Beispiel:} Hardware\authorA}

% Hier beschreibt Autor A seinen Teil der Implementierung.
% Welche Technologien verwendet wurden und welche Herausforderungen auftraten.

\subsection{Implementierungsdetails}

\emph{Hier werden die konkrete technische Umsetzung der Systemkomponenten sowie wichtige Designentscheidungen beschrieben. Code-Snippets, UML-Diagramme und Erklärungen zu kritischen Algorithmen verdeutlichen die Implementierung.}

\placeholderbox{Beispiel}{%
Die Kernkomponente analysiert Eingabedaten durch einen dreistufigen Prozess: (1) Daten-Validierung mit JSON-Schema, (2) Transformation mittels einer Custom-Pipeline mit Pluggable-Architektur, (3) Ausgabe-Serialisierung in mehreren Formaten. Die Daten-Validierung wurde als Middleware implementiert, um Fehler früh zu erkennen. Besondere Aufmerksamkeit galt der Fehlerbehandlung mit aussagekräftigen Fehlermeldungen für End-User sowie detaillierten Logs für Developer.%
    %
    \newline\newline
    ...
}

\subsection{Verwendete Technologien}

\emph{Dieser Abschnitt listet alle eingesetzten Programmiersprachen, Frameworks, Bibliotheken und Tools auf und begründet deren Auswahl. Die Darstellung erfolgt strukturiert mit konkreten Versionsangaben und Einsatzzwecken.}

\placeholderbox{Beispiel}{%
\textbf{Backend:} Node.js 18.14.0 mit Express.js 4.18.2 für schnelle API-Entwicklung; Mongoose 7.0.0 für die MongoDB-Integration. Die Wahl fiel auf Node.js wegen der Performance und dem großen Ökosystem.

\textbf{Frontend:} React 18.2.0 mit TypeScript 4.9 für Type-Safety; Material-UI 5.11 für konsistentes Design. React wurde gewählt für Component-Reusability und Virtual DOM Performance.%
    %
    \newline\newline
    ...
}


\section{\emph{Beispiel:} Frontend\authorA}

\subsection{Implementierungsdetails}

\ldots

\subsection{Verwendete Technologien}

\ldots
